\documentclass[11pt,a4paper]{article}
\usepackage[utf8]{inputenc}
\usepackage[T1]{fontenc}
\usepackage[french]{babel}
\usepackage[margin=2.4cm]{geometry}
\usepackage{amsmath,amssymb}
\usepackage{booktabs}
\usepackage{xcolor}
\usepackage[colorlinks=true,linkcolor=blue!50!black,citecolor=blue!50!black]{hyperref}
\usepackage{tcolorbox}
\tcbuselibrary{skins,breakable}
\usepackage{titlesec}
\usepackage{enumitem}
\usepackage{parskip}
\setlength{\parskip}{0.4em}

\definecolor{bleuprincip}{HTML}{1F4E78}
\definecolor{bleuclair}{HTML}{2E74B5}
\definecolor{orangealerte}{HTML}{C55A11}
\definecolor{vertok}{HTML}{548235}

\titleformat{\section}{\Large\bfseries\color{bleuprincip}}{\thesection}{0.6em}{}
\titleformat{\subsection}{\large\bfseries\color{bleuclair}}{\thesubsection}{0.6em}{}

\newtcolorbox{fichebox}[1][]{colback=vertok!6,colframe=vertok,boxrule=0.6pt,arc=2pt,
  left=7pt,right=7pt,top=5pt,bottom=5pt,breakable,#1}
\newtcolorbox{alertebox}[1][]{colback=orangealerte!7,colframe=orangealerte,boxrule=0.6pt,arc=2pt,
  left=7pt,right=7pt,top=5pt,bottom=5pt,breakable,title={\small\bfseries Limite de vérification, assumée},#1}
\newtcolorbox{noteboxbleu}[1][]{colback=bleuclair!7,colframe=bleuclair,boxrule=0.5pt,arc=2pt,
  left=6pt,right=6pt,top=4pt,bottom=4pt,breakable,title={\small\bfseries Note méthodologique},#1}

\title{\vspace{-1.3cm}\textcolor{bleuprincip}{\Huge\bfseries GDPNow-Maroc}\\[0.3cm]
\textcolor{black}{\Large Phase 2 --- Sélection des indicateurs}\\[0.15cm]
\textcolor{gray}{\large Justifications économiques et filtres --- Branche Électricité, gaz, eau}\\[0.1cm]
\textcolor{orangealerte}{\normalsize Sans contrainte de fraîcheur}}
\author{}
\date{\today}

\begin{document}
\maketitle
\thispagestyle{empty}

\begin{abstract}
\noindent Ce document présente le filtre économique a priori et le filtre de qualité de la donnée appliqués à la branche Électricité, gaz, eau. \textbf{81 indicateurs} sont retenus --- la base de candidats la plus large rencontrée dans cette procédure, avec une réserve méthodologique explicite sur l'exhaustivité de la vérification des doublons.
\end{abstract}

\tableofcontents
\vspace{0.3cm}

% ============================================================================
\section{Contexte --- une branche exceptionnellement documentée}
% ============================================================================

\begin{fichebox}
\textbf{128 indicateurs bruts} extraits, sur 165 colonnes sources --- la feuille la plus dense rencontrée dans cette procédure, reflétant le suivi très fin de la production électrique marocaine (par filière, par centrale nommée : JLEC, Tahaddart, Tarfaya, parcs éoliens individuels...). Après filtre économique et filtre de fréquence : \textbf{81 candidats}. Après filtre de qualité : \textbf{81 candidats retenus} --- aucun rejet.
\end{fichebox}

\begin{alertebox}
\textbf{Contrairement aux branches précédentes, la vérification systématique des doublons entre séries n'a pas été menée exhaustivement ici}, compte tenu du volume (128 indicateurs, dont 58 dans la seule catégorie « production détaillée »). Deux tables sources se recoupent largement en contenu (\textit{« Énergie appelée nette (mensuel).csv »} et \textit{« energie.csv »}), sans qu'il ait été vérifié, colonne par colonne, si certaines séries sont des duplicata exacts --- comme cela a été fait rigoureusement pour Agriculture, Pêche et Industrie de transformation. \textbf{Cette limite est assumée et signalée explicitement}, plutôt que de prétendre à une vérification qui n'a pas été menée à ce niveau de détail pour cette branche précise.
\end{alertebox}

% ============================================================================
\section{Le filtre économique a priori, par catégorie}
% ============================================================================

\subsection{Production d'électricité --- détail par filière et par centrale (58 séries retenues)}

\textbf{Source} : ONEE / Manar-Stat (tables \textit{« Énergie appelée nette »} et \textit{« energie.csv »}).

\textbf{Justification économique} : la production d'électricité par filière (hydraulique, thermique, éolienne, solaire) et par centrale nommée (JLEC, Tahaddart, Tarfaya, divers parcs éoliens) constitue l'extrant physique direct de la branche --- le lien le plus immédiat concevable avec sa valeur ajoutée.

\subsection{Ventes d'électricité par catégorie de clientèle (9 séries retenues)}

\textbf{Source} : ONEE, ventes mensuelles.

\textbf{Justification économique} : la ventilation par type de client (distributeurs, clients basse/moyenne/très haute tension) mesure la demande finale d'électricité, complémentaire à la production --- utile notamment si la VA de la branche suit davantage la commercialisation que la production brute.

\subsection{Consommation de combustibles par l'ONEE (5 séries retenues)}

\textbf{Source} : ONEE, consommation mensuelle en combustibles.

\textbf{Justification économique} : fioul, charbon (local et importé), gasoil, petcoke --- les intrants directs de la production thermique d'électricité, un indicateur avancé plausible de l'activité de production à venir.

\subsection{Crédit bancaire sectoriel (3 séries retenues)}

\textbf{Source} : Bank Al-Maghrib.

\textbf{Justification économique} : financement direct du secteur --- même raisonnement que pour les autres branches.

\subsection{Prix des matières premières énergétiques (2 séries retenues)}

\textbf{Source} : Platts / Bloomberg (via Manar-Stat).

\textbf{Justification économique} : le prix du Brent et du Butane, bien que fixés sur les marchés internationaux plutôt qu'au Maroc, influencent directement les coûts de production et les décisions d'investissement du secteur énergétique national --- un indicateur avancé de nature différente des précédents (un intrant de coût, pas un volume physique).

\subsection{Intrants de la raffinerie (2 séries retenues)}

\textbf{Source} : entrées de la raffinerie, mensuel.

\textbf{Justification économique} : pétrole brut et autres produits pétroliers entrant en raffinerie --- un indicateur avancé de l'activité de transformation énergétique à venir.

\subsection{Production et consommation agrégées (2 séries retenues)}

\textbf{Source} : BDD SECTORIEL (ONEE / Manar-Stat).

\textbf{Justification économique} : mesures agrégées de production d'électricité et de consommation de combustibles, complémentaires aux séries détaillées ci-dessus --- utiles comme indicateurs de synthèse si les séries granulaires s'avèrent trop bruitées individuellement.

% ============================================================================
\section{Deux séries écartées par cohérence avec le travail déjà mené}
% ============================================================================

\begin{noteboxbleu}
Deux séries annuelles présentes dans la base (\textit{« Importation du ciment hydraulique »}, \textit{« Consommation du sucre »}) auraient de toute façon été écartées par le filtre de fréquence --- mais leur exclusion est également cohérente avec un principe déjà établi dans les tout premiers travaux de ce projet : le ciment est un intrant de la Construction, pas de l'énergie ; le sucre n'a aucun lien mécanique identifiable avec cette branche.
\end{noteboxbleu}

% ============================================================================
\section{Synthèse}
% ============================================================================

\begin{fichebox}
\begin{center}
\begin{tabular}{lr}
\toprule
\textbf{Étape} & \textbf{Nombre de candidats} \\
\midrule
Indicateurs bruts extraits & 128 \\
Après filtre économique + fréquence & 81 \\
Après filtre de qualité (sans fraîcheur) & \textbf{81} \\
\bottomrule
\end{tabular}
\end{center}
\end{fichebox}

\begin{center}
\small
\begin{tabular}{lr}
\toprule
\textbf{Catégorie} & \textbf{Retenus} \\
\midrule
Production d'électricité --- détail & 58 \\
Ventes d'électricité par clientèle & 9 \\
Consommation de combustibles & 5 \\
Crédit bancaire sectoriel & 3 \\
Prix des matières premières & 2 \\
Intrants de la raffinerie & 2 \\
Production/consommation agrégées & 2 \\
\midrule
\textbf{Total} & \textbf{81} \\
\bottomrule
\end{tabular}
\end{center}

\end{document}
